\section*{Chapter \ref{ch:g9:elasticity-first-look} --- Elasticity: A First Look}

\begin{solution}{exo:g9:elasticity-first-look:1}
Most \omterm{def:g9:elasticity-first-look:elastic}{elastic} food: meals away from home ($0.81$). Least \omterm{def:g9:elasticity-first-look:elastic}{elastic} food: eggs ($0.27$). Lowest bar of the
chart: gasoline in the short run ($0.09$).
\end{solution}

\begin{solution}{exo:g9:elasticity-first-look:2}
$5 \div 10 = 0.5$.
\end{solution}

\begin{solution}{exo:g9:elasticity-first-look:3}
$0.5$: \omterm{def:g9:elasticity-first-look:inelastic}{inelastic} (below 1). $1.8$: \omterm{def:g9:elasticity-first-look:elastic}{elastic} (above 1).
\end{solution}

\begin{solution}{exo:g9:elasticity-first-look:4}
$3 \times 500 = 1\,500$ dalens; $4 \times 400 = 1\,600$ dalens.
\end{solution}

\begin{solution}{exo:g9:elasticity-first-look:5}
The quantity falls by $0.4 \times 10 = 4\,\%$. \omterm{def:g7:firms:revenue}{Revenue} is multiplied by $1.10 \times 0.96 = 1.056$: it
rises by $5.6\,\%$.
\end{solution}

\begin{solution}{exo:g9:elasticity-first-look:6}
The quantity falls by $1.5 \times 10 = 15\,\%$. \omterm{def:g7:firms:revenue}{Revenue} is multiplied by $1.10 \times 0.85 = 0.935$: it
falls by $6.5\,\%$.
\end{solution}

\begin{solution}{exo:g9:elasticity-first-look:7}
\omterm{def:g7:money:price}{Price}: $(5 - 4) \div 4 = 25\,\%$ rise. Quantity: $(800 - 700) \div 800 = 12.5\,\%$ fall. Elasticity
$12.5 \div 25 = 0.5$.
\end{solution}

\begin{solution}{exo:g9:elasticity-first-look:8}
First year: $0.09 \times 10 = 0.9\,\%$. After several years: $0.31 \times 10 = 3.1\,\%$.
\end{solution}

\begin{solution}{exo:g9:elasticity-first-look:9}
No. Other things did not stay the same: with the bridge closed, many people who usually drive took the
bus, so demand for bus rides rose. The fare rise alone would have lowered the number of rides; the
closed bridge raised it by more.
\end{solution}

\begin{solution}{exo:g9:elasticity-first-look:10}
At 4: $4 \times 800 = 3\,200$; at 5: $5 \times 700 = 3\,500$; at 6: $6 \times 600 = 3\,600$; at 7: $7
\times 500 = 3\,500$; at 8: $8 \times 400 = 3\,200$ dalens. The largest \omterm{def:g7:firms:revenue}{revenue} is at 6 dalens.
\end{solution}

\begin{solution}{exo:g9:elasticity-first-look:11}
First good: the quantity falls by $0.3 \times 10 = 3\,\%$, to 970 units: \omterm{def:g8:taxes-and-public-spending:tax}{tax} $970 \times 1 = 970$
dalens. Second good: the quantity falls by $2.0 \times 10 = 20\,\%$, to 800 units: \omterm{def:g8:taxes-and-public-spending:tax}{tax} 800 dalens. The \omterm{def:g8:taxes-and-public-spending:tax}{tax}
on the good with \omterm{def:g9:elasticity-first-look:inelastic}{inelastic demand} raises more.
\end{solution}

\begin{solution}{exo:g9:elasticity-first-look:12}
Other things may have changed between the two years (\omterm{def:g7:households:income}{incomes}, the \omterm{def:g7:money:price}{prices} of other goods, the number of
buyers), so the change in quantity is not due to the \omterm{def:g7:money:price}{price} alone. And the response two years apart may
be a short-run or a partly long-run response, which differ. (Also: the number depends on the point from
which percentage changes are measured.)
\end{solution}

\begin{solution}{pb:g9:elasticity-first-look:1}
\textbf{1.} $2.00 \times 10\,000 = 20\,000$ dalens; $2.50 \times 9\,000 = 22\,500$ dalens.
\textbf{2.} $(2.50 - 2.00) \div 2.00 = 25\,\%$.
\textbf{3.} $(10\,000 - 9\,000) \div 10\,000 = 10\,\%$.
\textbf{4.} It rose by $22\,500 - 20\,000 = 2\,500$ dalens.
\textbf{5.} $10 \div 25 = 0.4$.
\textbf{6.} \omterm{def:g9:elasticity-first-look:inelastic}{Inelastic}.
\textbf{7.} Riders fall by $1.5 \times 25 = 37.5\,\%$: $10\,000 \times 0.625 = 6\,250$ riders; \omterm{def:g7:firms:revenue}{revenue} $2.50
\times 6\,250 = 15\,625$ dalens.
\textbf{8.} With an elasticity of $0.4$, riders fall by a smaller percentage ($10\,\%$) than the fare rises
($25\,\%$), so \omterm{def:g7:firms:revenue}{revenue} rises; with $1.5$, they fall by a larger percentage ($37.5\,\%$), so \omterm{def:g7:firms:revenue}{revenue} falls.
\textbf{9.} The fare falls by $0.50 \div 2.50 = 20\,\%$; riders rise by $0.4 \times 20 = 8\,\%$: $9\,000 \times
1.08 = 9\,720$ riders. \omterm{def:g7:firms:revenue}{Revenue} $2.00 \times 9\,720 = 19\,440$ dalens. (Not back to $10\,000$: the
percentages are measured from a different starting point.)
\textbf{10.} $19\,440 - 22\,500 = -3\,060$ dalens a day.
\textbf{11.} Riders rise by $1.2 \times 20 = 24\,\%$: $9\,000 \times 1.24 = 11\,160$ riders. \omterm{def:g7:firms:revenue}{Revenue} $2.00
\times 11\,160 = 22\,320$ dalens.
\textbf{12.} $22\,320 - 22\,500 = -180$ dalens a day: in the long run, the fare cut costs the town almost
nothing.
\end{solution}
